\paragraph{Question simple (réponse confiante)}
\begin{lstlisting}
How do I load a model in 4-bit using bitsandbytes?
\end{lstlisting}
\begin{lstlisting}[backgroundcolor=\color{black!5}]
Answer (extracted from the documentation)

- The example below loads an 8B model in 4 bits with bitsandbytes. [1]
- To fine-tune a model that doesn't fit in memory, load it in 4-bit with bitsandbytes and train LoRA adapters with PEFT. [2]
- Make sure you have the latest bitsandbytes version so you can serialize 4-bit models and push them to the Hub with `PreTrainedModel.push_to_hub`. [3]

Code example [1]
```py
import torch
from transformers import AutoModelForCausalLM, BitsAndBytesConfig

quantization_config = BitsAndBytesConfig(
    load_in_4bit=True,
    bnb_4bit_compute_dtype=torch.bfloat16,
)
model = AutoModelForCausalLM.from_pre [...]
\end{lstlisting}
\noindent\textit{Décision : confident, confiance 88\%}.

\paragraph{Question de suivi (historique)}
\begin{lstlisting}
How do I enable gradient accumulation in the Trainer?
\end{lstlisting}
\begin{lstlisting}[backgroundcolor=\color{black!5}]
Answer (extracted from the documentation)

- This normalizes by tokens rather than a fixed step count with `gradient_accumulation_steps`. [1]
- Otherwise, `Trainer` divides loss by `gradient_accumulation_steps`. [1]
- If you subclass `Trainer` and override `Trainer.compute_loss`, set `loss_is_scaled_for_ga` to say how your loss is normalized: `False` for a per-batch mean, which `Trainer` then divides by `gradient_accumulation_steps`, or `True` if you already divided it by `num_items_in_batch`. [1]

Code example [2]
```py
from transformers import TrainingArguments

args = TrainingArguments(
    ...,
    per_device_train_batch_size=8,
  [...]
\end{lstlisting}
\noindent\textit{Décision : confident, confiance 68\%}.

\begin{lstlisting}
and how does it affect the loss?
\end{lstlisting}
\begin{lstlisting}[backgroundcolor=\color{black!5}]
Answer (extracted from the documentation)  I am not certain these excerpts answer your question exactly.

- For a custom loss function, include `num_items_in_batch` so `Trainer` divides the loss by the number of prediction targets across all mini-batches. [1]
- Otherwise, `Trainer` divides loss by `gradient_accumulation_steps`. [1]
- The loss shifts labels so the prediction at position `i` targets the token at position `i + 1`, which leaves position 0 of every sequence without a target. [1]

Code example [1]
```py
import torch.nn.functional as F

def compute_loss(outputs, labels, num_items_in_batch=None):
    logits = outputs["logits" [...]
\end{lstlisting}
\noindent\textit{Décision : uncertain, confiance 54\%, escalade proposée}.

\paragraph{Hors périmètre (refus net)}
\begin{lstlisting}
What is the tallest mountain in Africa?
\end{lstlisting}
\begin{lstlisting}[backgroundcolor=\color{black!5}]
I couldn't find a reliable answer to this in the Hugging Face Transformers documentation, so I would rather not guess. I can pass your question to a human expert.
\end{lstlisting}
\noindent\textit{Décision : no\_answer, confiance 5\%, escalade proposée}.

\paragraph{Hors périmètre proche du domaine (réponse avec réserve : cas limite)}
\begin{lstlisting}
How do I deploy my web app on AWS Lambda?
\end{lstlisting}
\begin{lstlisting}[backgroundcolor=\color{black!5}]
Answer (extracted from the documentation)  I am not certain these excerpts answer your question exactly.

- Add launch to create a web server and start up the app. [1]
- The web app runs on a local server by default. [1]
- Deploy the model locally from the command line with llama-cli or start a web UI with llama-server. [3]

Code example [2]
```py
from starlette.applications import Starlette
from starlette.responses import JSONResponse
from starlette.routing import Route
from transformers import pipeline
import asyncio

async def homepage(request):
    payload = await request.body()
    string = payload.decode("utf-8")
    response_q  [...]
\end{lstlisting}
\noindent\textit{Décision : uncertain, confiance 22\%, escalade proposée}.

\paragraph{Question vague (incertitude / escalade)}
\begin{lstlisting}
My model is slow.
\end{lstlisting}
\begin{lstlisting}[backgroundcolor=\color{black!5}]
I couldn't find a reliable answer to this in the Hugging Face Transformers documentation, so I would rather not guess. I can pass your question to a human expert.
\end{lstlisting}
\noindent\textit{Décision : no\_answer, confiance 16\%, escalade proposée}.

\paragraph{Demande d'un humain}
\begin{lstlisting}
Can I speak to a human please?
\end{lstlisting}
\begin{lstlisting}[backgroundcolor=\color{black!5}]
Understood  I'm creating a ticket for a human expert, with a summary of our conversation.
\end{lstlisting}
\noindent\textit{Décision : human\_requested, confiance 100\%, escalade proposée}.
